\begin{annexitem}{annex:p12.har.hurkens-beta}{thm:p12-har-hurkens-beta}{All universe/recipe typings and the actual beta identity}
\begin{proof}
Work first under X:Box. Both X and the sort star have type Box,
so Box/Box forms \(\wp X=X\to\ast:\Box\), then
\(\wp\wp X:\Box\), then
\((\wp\wp X\to X)\to\wp\wp X:\Box\).
The domain annotation Box has type triangle; the admitted triangle/Box
pair therefore forms the displayed \(\mathcal U:\Box\).
Box/Box also gives \(\wp\mathcal U:\Box\) and
\(\wp\wp\mathcal U:\Box\). Falsehood is well formed at star:
under p:star the variable p has type star, and Box/star forms
\(\Pi p:\ast.p:\ast\). Implications between small formulas
use star/star, with fresh unused proof binders.

For the tau recipe take \(t:\wp\wp\mathcal U\),
\(X:\Box\), \(f:\wp\wp X\to X\), \(p:\wp X\),
and a fresh \(x:\mathcal U\). Universe application gives
\(\{xX\}:(\wp\wp X\to X)\to\wp\wp X\);
application to f gives \(\wp\wp X\); applying f again gives
X; applying p gives a term of type star. Hence
\(\lambda x:\mathcal U.p(f(\{xX\}f)):\wp\mathcal U\).
Application of t to this argument also has type star. Abstraction over
p gives a term of type \(\wp\wp X\), abstraction over f gives
the stated fiber of U, and polymorphic abstraction over X gives
\(\tau t:\mathcal U\). Every product in this chain is
Box/Box except the final triangle/Box universe product.

For \(s:\mathcal U\), instantiate its universe parameter at
the actually formed \(\mathcal U:\Box\). The result has type
\((\wp\wp\mathcal U\to\mathcal U)\to\wp\wp\mathcal U\).
The already checked tau instances give
\(\lambda t:\wp\wp\mathcal U.\tau t:
\wp\wp\mathcal U\to\mathcal U\). Application therefore
gives \(\sigma s:\wp\wp\mathcal U\). In particular
\(\tau\sigma y:\mathcal U\), and
\(q_p=\lambda y:\mathcal U.p(\tau\sigma y):\wp\mathcal U\).
This validates the types before the computational use.

Let \(F=\lambda u:\wp\wp\mathcal U.\tau u\), a notation
for this actual checked expression. Expand sigma and the tau instance:
\[
\begin{aligned}
 (\sigma(\tau t))p
 &=((\tau t)\mathcal U)F p\\
 &\twoheadrightarrow_\beta
 t(\lambda x:\mathcal U.p(F(\{x\mathcal U\}F)))\\
 &\twoheadrightarrow_\beta
 t(\lambda x:\mathcal U.p(\tau(\{x\mathcal U\}F)))\\
 &=t(\lambda x:\mathcal U.p(\tau\sigma x))=tq_p.
\end{aligned}
\]
The equal signs on the first/last lines expand declared abbreviations;
the intervening steps are actual beta contractions of the displayed
annotated lambdas. No equation constant, extensionality or fresh logical
axiom has been supplied. Specializing t to sigma x consequently gives
\(\sigma(\tau\sigma x)p=_\beta\sigma x q_p\), the identity
needed in both proof branches below.
\end{proof}
\end{annexitem}

\begin{annexitem}{annex:p12.har.hurkens-bottom}{thm:p12-har-hurkens-bottom}{Every small-formula formation and the R0/M0/L0 application chain}
\begin{proof}
For \(p:\wp\mathcal U\), both \(\sigma x p\) and px
have type star. Their implication has type star by star/star; quantifying
the data variable x:U uses Box/star and gives \(I(p):\ast\).
The same argument gives \(K(x):\ast\), since \(\tau\sigma x:U\).
Negation \(K(x)\Rightarrow\bot\) is small, so abstraction over
x gives \(\Delta:\wp\mathcal U\). Abstraction of I over
p gives \(\lambda p:\wp\mathcal U.I(p):\wp\wp\mathcal U\),
whose tau instance has type U. Expanding this one constructed instance
and beta-contracting its application to \(q_{X,f,p}\) yields exactly
the displayed finite Omega expression. Its further q application is a
variable substitution, so the indicated normal form exists by this
finite calculation; global normalization is never assumed.

Thus \(\Omega:\mathcal U\), and the small formula W is formed
by star/star implication inside Box/star quantification over p.
The computed identity, at the actual I expression, gives
\[
 \sigma\Omega\,p=_\beta I(q_p).
\]
Also \(q_p x=_\beta p(\tau\sigma x)\). These are beta
equalities between well-typed small formulas; the source conversion rule
may use them without a separate proof of equality.

Check R0 in its complete context \(p:\wp\mathcal U,1:I(p)\).
Instantiation at Omega gives
\(1\Omega:\sigma\Omega p\Rightarrow p\Omega\).
For fresh x:U, instantiation at tau sigma x gives
\(1(\tau\sigma x):
\sigma(\tau\sigma x)p\Rightarrow p(\tau\sigma x)\).
By the first lemma its domain converts to \(\sigma xq_p\),
and its codomain is \(q_p x\). Hence
\(\lambda x:\mathcal U.1(\tau\sigma x):I(q_p)\), which
converts to \(\sigma\Omega p\). Application to this actual
proof yields \(p\Omega\). Abstracting over proof label1 and
over p gives exactly \(R_0:W\).

Check M0 under \(x:\mathcal U,2:\sigma x\Delta,3:K(x)\).
Instantiation at Delta gives
\(3\Delta:\sigma x\Delta\Rightarrow\Delta(\tau\sigma x)\),
so \((3\Delta)2:\Delta(\tau\sigma x)
=K(\tau\sigma x)\Rightarrow\bot\).
For fresh p:PU, the actual instance \(3q_p\) has type
\(\sigma xq_p\Rightarrow q_p(\tau\sigma x)\).
The domain converts to \(\sigma(\tau\sigma x)p\) by the
computed identity, and the codomain converts to
\(p(\tau\sigma(\tau\sigma x))\). Therefore
\(\lambda p:\wp\mathcal U.3q_p:K(\tau\sigma x)\).
Applying the previously obtained negation to this proof gives bottom.
Abstraction over label3 gives Delta x; abstraction over label2 gives
\(\sigma x\Delta\Rightarrow\Delta x\); abstraction over x
gives precisely \(M_0:I(\Delta)\). Every premise has been supplied,
rather than assuming that Delta is inductive.

Check L0 under its sole proof assumption \(0:W\).
Instantiation at Delta gives
\(0\Delta:I(\Delta)\Rightarrow\Delta\Omega\),
so \((0\Delta)M_0:\Delta\Omega=K(\Omega)\Rightarrow\bot\).
For fresh p:PU, \(0q_p:I(q_p)\Rightarrow q_p\Omega\).
Its domain converts to sigma Omega p, and its codomain to
\(p(\tau\sigma\Omega)\), so
\(\lambda p:\wp\mathcal U.0q_p:K(\Omega)\).
The final application yields bottom. Abstraction over label0 gives
\(L_0:W\Rightarrow\bot\). The previously constructed R0 supplies
W in the empty context; thus \(D_0=L_0R_0:\bot\) is closed.
All helper types used only the exact four source pairs. In particular
no triangle/star quantification, ex-falso elimination, double-negation
axiom or primitive sigma/tau equality was inserted.
\end{proof}
\end{annexitem}

\begin{annexitem}{annex:p12.har.star-inconsistency}{thm:p12-har-star-inconsistency}{The actual sort collapse and the resulting inhabitant}
\begin{proof}
The constant sort map sends both source axioms to the one admitted
star:star axiom, and every one of the four source product triples
to star/star/star. It fixes variables, so freshness remains valid,
and commutes with capture-avoiding substitution by induction on terms.
The seven-rule transport of \cref{thm:p12-har-pts-morphism} therefore
maps the complete closed derivation for D0 to a closed lambda-star
derivation for its explicitly expanded image Dstar. The false-type
expression \(\Pi p:\ast.p\) has that same image; consequently
\(\vdash D_\ast:\Pi p:\ast.p\).

For an arbitrary valid target context, start from this closed derivation
and its actual prefix certificates. Such certificates follow inductively
from occurrence in a derivable judgment: sort uses the empty context,
variable/weakening exhibit their extension annotation, and the other
rules inherit their base context from a premise.
Weaken through the context's actual annotation certificates, one
fresh declaration at a time. This gives Dstar the same type in Gamma.
Given the actual premise \(\Gamma\vdash A:\ast\), product
elimination then derives \(\Gamma\vdash D_\ast A:A\).
No inhabitant of A was used to construct the context, map or argument;
the paradox term supplies it. The source hierarchy has been transported,
not silently identified with lambda-star or any other universe system.
\end{proof}
\end{annexitem}

\begin{annexitem}{annex:p12.har.star-no-normal-form}{thm:p12-har-star-no-normal-form}{Typed beta preservation and the actual normal-form obstruction}
\begin{proof}
First supply typed beta preservation for this target rather than importing
the inconsistent conclusion as a premise. Valid-context projection and
thinning follow by the variable/weakening induction. Generation strips
final conversion/weakening to expose the subject's constructor premises;
raw confluence gives product injectivity. Type correctness here has no
untyped-top exception: every derivable type is itself typed at star.
For sort/product this uses star:star; variables use their annotation;
weakening and conversion retain the typed result; abstraction has its
product-formation premise; application obtains the codomain sort by
generation of its typed function product and substitutes the actual
argument by \cref{thm:p12-har-pts-substitution}.

Context conversion replaces a typed binder annotation by a typed
beta-equal one: derive the selected variable at the new annotation,
convert back to the old annotation using its existing sort judgment,
then transform each later annotation before its variable/weakening step.
Product/lambda/application/conversion rebuild their same target rule.
For subject reduction induct on the subject. Annotation contraction
uses the smaller annotation typing, this context conversion, and
rebuilds the binder at the same single admitted triple. A body or
application-child contraction uses its smaller-subject induction.
At a root lambda application, generation and product injectivity
convert the argument to its literal annotation; whole-context
substitution types the contractum, and type correctness supplies the
result-type sort for restoring its generated and original types.
These are all compatible beta positions. Hence every finite reduct
of Dstar still has type \(\Pi p:\ast.p\).

Now there is no closed beta-normal term of that type. A closed normal
term is a lambda, a product, the sort star, or an application spine.
Any product or sort has generated type star; that normal sort cannot
convert to a product type because raw confluence preserves the outer
product constructor. An application spine cannot start with a lambda
(that would be a redex); a closed spine cannot start with a free
variable; and a product or star head has generated type star, not a
function product, so it cannot supply the application premises. Thus
a putative normal inhabitant must be \(\lambda p:A.N\).

Its generated product type converts to \(\Pi p:\ast.p\).
Product injectivity gives \(A=_\beta\ast\) and its body type
beta-equal to p. Its annotation is normal, so A is literally star
up to irrelevant binder names. In the context \(p:\ast\),
conversion consequently gives \(N:p\); N is also normal.
But a normal term in this context has only p as a possible free head.
That variable has generated type star, not p, and those are distinct
normal forms. Sort/product heads likewise have type star; lambdas
have product types; none converts to p. An application spine again
has no function-typed head. No such N exists.

If Dstar had a normal form, typed preservation would give just such a
closed normal inhabitant, a contradiction. Hence this explicit typed
term has no normal form. Equivalently, always choosing the first beta
redex in a fixed finite traversal gives an infinite path: a finite
termination point would be a normal form. This proves failure of weak
normalization, hence of strong normalization, for the exact target.
No claim about every stratified universe calculus follows.
\end{proof}
\end{annexitem}
